\documentclass{article}

\usepackage{amsmath,amssymb}
\usepackage{xurl}
\usepackage[hidelinks]{hyperref}
\setlength{\emergencystretch}{2em}

% Shared commands for notes in docs/paper-gaps/.

\DeclareUnicodeCharacter{2080}{\ifmmode{}_{0}\else\textsubscript{0}\fi}
\DeclareUnicodeCharacter{2081}{\ifmmode{}_{1}\else\textsubscript{1}\fi}
\DeclareUnicodeCharacter{2082}{\ifmmode{}_{2}\else\textsubscript{2}\fi}
\DeclareUnicodeCharacter{2099}{\ifmmode{}_{n}\else\textsubscript{n}\fi}

\newcommand{\Lean}{\textsc{Lean}}
\ExplSyntaxOn
\NewDocumentCommand{\leanid}{m}{
  \ifmmode
    \textsf{\detokenize{#1}}
  \else
    \group_begin:
    \urlstyle{sf}
    \tl_set:Nn \l_tmpa_tl {#1}
    \tl_replace_all:Nnn \l_tmpa_tl {\_} {_}
    \exp_args:NV \path \l_tmpa_tl
    \group_end:
  \fi
}
\ExplSyntaxOff
\providecommand{\tr}{\operatorname{tr}}
\newcommand{\tnleanIssueBase}{https://github.com/LionSR/TNLean/issues/}
\newcommand{\tnleanPrBase}{https://github.com/LionSR/TNLean/pull/}
\newcommand{\tnleanDocBase}{https://sirui-lu.com/TNLean/docs/}
\newcommand{\ghissue}[1]{\href{\tnleanIssueBase#1}{GitHub issue \##1}}
\newcommand{\ghpr}[1]{\href{\tnleanPrBase#1}{PR \##1}}
\newcommand{\doclink}[1]{\href{\tnleanDocBase#1.html}{\path{#1}}}

% Dirac notation and the identity operator, as in blueprint/src/macros/common.tex.
% \providecommand keeps note-local definitions authoritative.
\usepackage{dsfont}
\providecommand{\ket}[1]{|#1\rangle}
\providecommand{\bra}[1]{\langle #1|}
\providecommand{\braket}[2]{\langle #1 | #2 \rangle}
\providecommand{\ketbra}[2]{|#1\rangle\!\langle #2|}
\providecommand{\Id}{\mathds{1}}
\providecommand{\one}{\mathds{1}}
\providecommand{\C}{\mathbb{C}}
\providecommand{\MN}[1]{M_{#1}(\C)}
\providecommand{\MD}{M_D(\C)}

% External referencing for paper sources.  The build helper writes sanitized
% auxiliary files under build/paper-aux/ and excludes duplicated source labels.
\usepackage{xr}

% Import a generated paper auxiliary file with a caller-supplied label prefix.
% The candidate paths support builds from docs/paper-gaps/, the repository root,
% and build/paper-aux/ itself.
\newcommand{\importpaperaux}[2]{%
  \IfFileExists{../../build/paper-aux/#2.aux}{%
    \externaldocument[#1]{../../build/paper-aux/#2}%
  }{%
    \IfFileExists{build/paper-aux/#2.aux}{%
      \externaldocument[#1]{build/paper-aux/#2}%
    }{%
      \IfFileExists{#2.aux}{%
        \externaldocument[#1]{#2}%
      }{}%
    }%
  }%
}

% General external reference macros
\newcommand{\paperref}[2]{\ref{#1-#2}}
\newcommand{\papereqref}[2]{(\ref{#1-#2})}

% CPSV16 source (1606.00608 / MPDO-22-12-17-2.tex) external document and shortcuts
\importpaperaux{cpsv16-}{1606.00608/MPDO-22-12-17-2-tnlean}
\newcommand{\cpsvref}[1]{\paperref{cpsv16}{#1}}
\newcommand{\cpsveqref}[1]{\papereqref{cpsv16}{#1}}

% Verdict marker: \gapnote{<kind>}{<status>}. Kind is the policy
% classification (clarification, local-correction, scope-restriction,
% unfaithful, false-source, open-gap); status is open, wip (elimination
% underway), resolved, or historical. Machine-read by the site index;
% prints a verdict line.
\newcommand{\gapnote}[2]{\par\noindent\textbf{Verdict:} #1 (#2).\par}


\title{Proof-Route Alignment for the Periodic Overlap Dichotomy
  (de las Cuevas--Cirac--Schuch--P\'erez-Garc\'ia)}
\author{}
\date{2026-06-04; revised 2026-08-02}

\begin{document}
\maketitle

\section{Scope}

This note records, in mathematical terms, where the formal development of the
periodic overlap dichotomy chooses a proof route that differs from the appendix
argument of \cite[Appendix~A]{DeLasCuevas2017Irreducible}, and where its
statements restrict the source.\footnote{Formal files:
\leanid{TNLean/MPS/Periodic/Overlap/}. Local source:
\leanid{Papers/1708.00029/main.tex}. Paper-level umbrella:
\href{https://github.com/LionSR/TNLean/issues/619}{issue \#619}; proposition-level
subtracking: \href{https://github.com/LionSR/TNLean/issues/81}{issue \#81};
the Case-3 contraction and phase construction was tracked by
\href{https://github.com/LionSR/TNLean/issues/873}{issue \#873} and is now
formalized by \leanid{sectorTensor_proportional_of_blockedMatch}.}

The source proposition is \emph{equal-or-orthogonal-generalized}.\footnote{Stated
in \cite{DeLasCuevas2017Irreducible}, lines 589--609, and proved in Appendix~A,
lines 903--1118.} None of the deviations below is unfaithful in the
smuggled-hypothesis sense: each formal statement is a faithful (or explicitly
restricted) version of the source, and the route substitutions are
mathematically equivalent.

\section{Positive chain lengths}

The matrix-product-vector equalities used by the periodic overlap argument are
now stated for $N\geq 1$.  In particular,
\leanid{peripheralProportionalCase_periodicFT_of_sameMPV₂Pos} assumes equality
at every positive length, and
\leanid{PeripheralProportionalCaseRootFromRescaling} concludes equality at
every positive length after phase rescaling.  This is sufficient because the
overlap argument is asymptotic and uses only a tail of positive lengths.  The
empty word carries only a bond-dimension trace and plays no role in
\cite[theorem \texttt{thm:bd}, lines~613--623]{DeLasCuevas2017Irreducible}.
Likewise, propagation of nonzero cyclic-sector dimensions is derived directly
from the corner-algebra equivalences and the cyclic projector shift, rather
than from the empty-word trace identity.

\section{Different-period decay via the peripheral spectrum}

\textbf{Source.}\footnote{\cite{DeLasCuevas2017Irreducible}, Appendix~A,
lines 917--950.} For periods $m_a\neq m_b$, the paper blocks at
$p=\operatorname{lcm}(m_a,m_b)$ and uses Lemma~\emph{bdcf}: the blocks
$P_uA^{(p)}$ and $Q_vB^{(p)}$ are non-repeated normal tensors. Applying the
one-site translation operator to $Q_vB^{(p)}$ versus $Q_{v+1}B^{(p)}$ yields a
contradiction unless $\lim_N|\langle V_N(A)|V_N(B)\rangle|=0$.

\textbf{\Lean{} route.}\footnote{Formal declaration:
\leanid{periodicOverlap_tendsto_zero_of_ne_period} in \leanid{DifferentPeriod.lean}.}
The same obstruction is realized spectrally. For a tensor $A$, write
$\mathcal E_A$ and $\mathcal E_A^*$ for the transfer map and its adjoint,
defined by
\begin{align}
    \mathcal E_A(Y) &= \sum_i A^iYA^{i\dagger},&
    \mathcal E_A^*(Y) &= \sum_i A^{i\dagger}YA^i .
\end{align}
If $D_a=D_b$ and $B^i=\zeta\,XA^iX^{-1}$, the comparison is
\begin{align}
    \mathcal E_B(Y)
        &= \zeta\bar\zeta\,
           X\mathcal E_A(X^{-1}YX^{-\dagger})X^\dagger .
    \tag{Case 1}
\end{align}
Eigenvalue uniqueness for an irreducible completely positive map
(\cite{Wolf2012Quantum}, Theorem~6.3) forces $|\zeta|=1$, so $\mathcal E_A$ and
$\mathcal E_B$ have equal peripheral spectra. An $m$-periodic tensor has
peripheral spectrum $\{e^{2\pi i r/m}:0\le r<m\}$, so equality forces $m_a=m_b$.
Thus different periods exclude gauge-phase equivalence, and
\emph{equal-or-orthogonal} (the irreducible trace-preserving overlap dichotomy)
gives the decay. This substitutes the paper's lcm-blocking + translation
argument; it is complete and proved.

\section{Sector non-repetition via the blocked fixed-point structure}

\textbf{Source.}\footnote{Lemma~\emph{bdcf} in
\cite{DeLasCuevas2017Irreducible}, lines 404--423.} For a periodic block,
$\mathcal E_A$ is irreducible with peripheral spectrum $\{\omega^r\}$,
$\omega=e^{2\pi i/m}$. The blocked map $\mathcal E_A^m$ then has $1$ as its only
modulus-one eigenvalue (multiplicity $m$), with fixed-point set exactly
$\{P_u\Lambda_AP_u\}_u$, while the adjoint $\mathcal E_A^{*m}$ has fixed points
$\{P_u\}_u$. If, for $u\neq v$, a gauge-phase equivalence
$C_u^{\mathbf i}=e^{i\xi}UC_v^{\mathbf i}U^\dagger$ held with $U=P_uUP_v$
($UU^\dagger=P_u$, $U^\dagger U=P_v$), then
\begin{align}
    \mathcal E_A^m(U)
        &= \sum_{\mathbf i} C_u^{\mathbf i}UC_v^{\mathbf i\dagger} \\
        &= e^{i\xi}U
           \sum_{\mathbf i} C_v^{\mathbf i}C_v^{\mathbf i\dagger} \\
        &= e^{i\xi}U .
        \tag{$\mathcal E_{C_v}(P_v)=P_v$}
\end{align}
so $U$ is a modulus-one eigenvector of $\mathcal E_A^m$. As the only such
eigenvalue is $1$ with the diagonal fixed points $\{P_w\Lambda_AP_w\}$ and
$U=P_uUP_v$ is off-diagonal for $u\neq v$, this is a contradiction.

\textbf{\Lean{} status.}\footnote{Formal declaration:
\leanid{not_gaugePhaseEquiv_of_orthogonal_cyclicSector_traces} in
\leanid{SelfOverlap.lean}.} The lemma statement is the faithful non-repetition fact:
distinct cyclic sectors (orthogonal projectors $P_uP_v=0$) are not gauge-phase
equivalent. The formal proof now follows the spectral route above.  From a
gauge-phase equivalence between two compressed sectors, it constructs a nonzero
off-diagonal eigenvector of $\mathcal E_A^m$, using the sector corner-letter
identity and the rectangular support isometries for the compressed corners.  The
existing off-diagonal vanishing lemma then gives the contradiction.  The scalar
normalization reuses the same Perron--Frobenius eigenvalue-uniqueness input as
the different-period proof.\footnote{Formal declarations:
\leanid{exists_offDiag_eigenvector_of_gaugePhase_cast_left},
\leanid{offDiag_eigenvector_eq_zero_of_isPeriodic}, and
\leanid{period_eq_of_gaugePhaseEquiv_of_isPeriodic}.}

\textbf{Signature correction (periodicity is load-bearing).} An adversarial
proof-attempt found that the lemma was \emph{underspecified}: with only the bare
cyclic-projection data (orthogonal $P_k$ summing to $1$, the adjoint shift
$\mathcal E_A^{*}(P_{k+1})=P_k$, blocked-letter commutation, the trace formula,
and orthogonality), the conclusion is \emph{false} --- a reducible $A=B\oplus B$
(for an irreducible period-$m$ block $B$) admits the same cyclic-projection data
yet has gauge-phase-equivalent distinct sectors. The spectral argument genuinely
requires $A$ to be a periodic (irreducible) block: that is what makes $\mathcal
E_A$ irreducible with the stated peripheral spectrum and the sectors primitive.
The hypothesis \leanid{hP : IsPeriodic m A} was therefore added to the lemma
signature (it is supplied at the unique call site
\leanid{sectorBlocks_not_gaugePhaseEquiv_of_ne}). With this hypothesis and the
corner-isometry data, the spectral proof is now formalized.

\section{Sector-match case: propagation and the phase construction}

\textbf{Source.}\footnote{\cite{DeLasCuevas2017Irreducible}, Appendix~A,
lines 961--1117; the initial blocked equation \emph{eq:Nm} appears at
lines 978--984.} Given one matched pair, applying the
translation operator $T^l$ to the blocked equation \emph{eq:Nm}
and invoking Theorem~2.10 of \cite{CiracPerezGarcia2017Matrix} produces, at each
offset, a phase and a unitary $U_{\tilde v+l}=P_{\tilde u+l}U_{\tilde
v+l}Q_{\tilde v+l}$, so the offset
$v-u=q$ is constant. Then each blocked product
$F_u$ becomes injective after $N_0$ repetitions, with a right
inverse $\Omega_u$; contracting the concatenated products
with the $\Omega_u$ recovers per-site proportionality $A_u^i=\kappa_v\,
e^{i\eta/m}\,B_v^i$. The phase
bookkeeping is load-bearing: $\prod_v\kappa_v=1$ (\emph{eq:prodkappaprop}) and
$|\kappa_v|=1$ from $\|\sum_iA_u^{i\dagger}A_u^i\|=1$, so $\kappa_v=e^{i\theta_v}$
with $\sum_v\theta_v=0$; choosing $\phi_v$ with $\theta_v=\phi_v-\phi_{v+1}$
telescopes the per-sector phases into a single global phase $\xi=\eta/m$ and a
single global unitary $U=\sum_u e^{i\phi_{u+q}}P_uU_{u+q}Q_{u+q}$, giving
$A^i=e^{i\xi}UB^iU^\dagger$.

\textbf{\Lean{} status.}\footnote{Formal declarations:
\leanid{sectorMatch_propagation} in
\leanid{SectorMatch/Propagation.lean}, and
\leanid{sectorTensor_proportional_of_blockedMatch} in
\leanid{SectorMatch/Contraction.lean}, and
\leanid{periodicOverlap_gaugeEquiv_of_sector_match} in
\leanid{SectorMatch/Consequences.lean}.}
The statements track this two-stage structure: the $T^l$ + Theorem~2.10
staircase, followed by the $\Omega_u$ contraction and the
$\kappa/\theta/\phi$ phase construction. The top-level theorem now invokes these
component lemmas; the contraction leaf
\leanid{sectorTensor_proportional_of_blockedMatch} and the resulting
unitary global gauge are proved, and the repeated-blocks theorem follows. The
docstrings have been realigned to cite the source equation labels and line
ranges (replacing earlier published-version labels such as ``A.8'' and
``A.14--A.18'' that do not occur in the local source), and the
$\kappa/\theta/\phi$ phase construction is now stated explicitly as
load-bearing.

\textbf{Formalized single-site inputs.} The single-site off-diagonal
decomposition is now formalized from the cyclic adjoint-shift
relation.\footnote{The formal declarations are
\leanid{offDiag_shift_of_adjoint_cyclic_shift},
\leanid{eq_sum_offDiag_of_adjoint_cyclic_shift},
\leanid{IsCyclicSectorDecomp.offDiag_shift}, and
\leanid{IsCyclicSectorDecomp.eq_sum_offDiag}.} It gives
$P_{u+1}A^i=A^iP_u$ and $A^i=\sum_u P_{u+1}A^iP_u$ for a chosen cyclic
sector decomposition. This is the inverse-indexed form of
\emph{eq:Aoffdiag}, lines~335--339; the later \emph{eq:Auprop}, line~1014,
defines the sector letters $A_u^i=P_uA^iP_{u+1}$.

The construction of the cyclic sectors also now records the letter-level
corner identity and a rectangular support isometry for each corner.\footnote{Formal
declarations:
\leanid{exists_cyclic_sector_decomp_with_letter_and_isometry_after_blocking_of_isPeriodic}
and its projected wrapper
\leanid{exists_cyclic_sector_decomp_with_letter_after_blocking_of_isPeriodic}.}
If $\varphi_u$ is the compression isomorphism from the compressed sector algebra
onto $P_uM_DP_u$, then
\[
    \varphi_u(C_u^{\mathbf i})=
        P_u (A^{(m)})^{\mathbf i} P_u .
\]
This is the exposed formal version of the blocked corner-letter relation in
Lemma~\emph{bdcf} and \emph{Cu}. Together with the support isometry, it transports
the compressed mixed-transfer eigenvector to the ambient corner and supplies the
off-diagonal modulus-one eigenvector used in the spectral contradiction.

Let $O_{C,D}(N):=\langle V_N(C)\mid V_N(D)\rangle$. The one-site
sector-overlap translation is also formalized.\footnote{Formal declaration:
\leanid{sectorOverlap_succ_eq_of_cyclicSectorDecomp}.} It proves, for every
$N>0$,
\[
    O_{A_{u+1},B_{v+1}}(N)=O_{A_u,B_v}(N).
\]
It is the overlap-sum form of the translation step in lines~985--1002, using
the single-site relations inside the paired trace sum and the repository's
inverse cyclic indexing convention.

\textbf{Formalized Case-3 contraction.} After these single-site inputs, the
contraction and gauge construction proceeds as follows:
\begin{itemize}
  \item The $m$-factor cyclic $\Omega_u$-contraction with the $\eta$ bookkeeping:
    starting from the common right inverses for the injective tensors $F_u$,
    contract the cyclic product in lines~1041--1068 to obtain the uniform
    product-tensor identity \emph{eq:resultprop}. The displayed inverse list
    in lines~1049--1052 has only $m-1$ factors, but the displayed
    concatenation, the $(mN_0+1)$ repetition count, and the definition of
    $\eta_v$ require one inverse for each of the $m$ inserted $F$-factors.
  \item The final normalization and gauge construction: use the proved scalar
    extraction theorem to obtain $\kappa_v$ with $\prod_v\kappa_v=1$, prove
    $|\kappa_v|=1$ from the left-canonical normalization, apply the finite-cycle
    coboundary lemma to write $\kappa_v=e^{i(\phi_v-\phi_{v+1})}$, and assemble
    the global unitary
    $U=\sum_u e^{i\phi_{u+q}}P_uU_{u+q}Q_{u+q}$.
\end{itemize}
These steps discharge the contraction obligation for the repeated-blocks
theorem.\footnote{Formal declaration:
\leanid{sectorTensor_proportional_of_blockedMatch} in
\leanid{SectorMatch/Contraction.lean}.} The tensor-product scalar
extraction and product-one normalization have already been isolated in
\leanid{PiTensorProductPhase.exists_kappa_of_piTensorProduct_eq_smul} and
\leanid{PiTensorProductPhase.exists_kappa_product_one_of_piTensorProduct_eq_root_smul};
the finite-cycle phase choice is isolated in
\leanid{exists_fin_complex_unit_cyclic_coboundary_shift_of_prod_eq_one}. The
one-step transport theorem is closed by the single-site overlap theorem above
together with the overlap-to-gauge theorem and the proved sector primitivity.\footnote{
The formal declarations are
\leanid{sectorGaugePhaseEquiv_succ_of_cyclicTransport} and
\leanid{gaugePhaseEquiv_of_overlap_norm_tendsto_one_of_irreducible_TP}.}

\textbf{Corner-letter building blocks.} The corner-transition tensors
$A_u^i=P_uA^iP_{u+1}$ (eq:Auprop) and their repeated products
$F_u^{\mathbf i}=A_u^{i_1}A_{u+1}^{i_2}\cdots$ (eq:Fu) are now defined as
\leanid{MPSTensor.cornerLetter} and \leanid{MPSTensor.cornerProd}. The cyclic
concatenation law $F_u^{\mathbf{i}_1\mathbf{i}_2}
=F_u^{\mathbf i_1}\,F_{u+\ell_1}^{\mathbf i_2}$, by which the appendix repeats a
single block around the cycle to an injective length, is
\leanid{MPSTensor.cornerProd_append}; the junction projector is absorbed by
idempotency, so the law needs no injectivity or normality input. These are the
$A_u^i/F_u$ building blocks that enter the first bullet above. The formal
right-inverse contraction of $F_u$ to \emph{eq:resultprop} and the unitary
global gauge assembly are completed in
\leanid{sectorTensor_proportional_of_blockedMatch}.

\section{Scope restriction: independence without spanning}

\textbf{Source.}\footnote{\cite{DeLasCuevas2017Irreducible}, lines 604--608
and Lemma~\emph{Lem1t}, lines 511--519. The phrase about the name
``basis of periodic vectors'' appears at line 611.} As a consequence of the
proposition and Lemma~\emph{Lem1t}, for $N\ge N_0$ the non-zero members of
$\{|V_N(A_j)\rangle\}_{j\in J}$ are linearly independent \emph{and} span
$|V_N(A)\rangle$; the spanning half is what ``justifies the name basis of
periodic vectors''.

\textbf{\Lean{} statement.}\footnote{Formal declaration:
\leanid{periodicBasis_eventuallyLinearlyIndependent} in \leanid{Dichotomy.lean};
the matching blueprint entry is in
\path{blueprint/src/chapter/ch22_periodic_ft_overlap_sector_match_and_consequences.tex}.}
Only the independence half is stated; the spanning clause is dropped. The
basis condition is encoded directly as the pairwise non-repetition hypothesis
\leanid{hNonrep} together with
a common-period divisibility \leanid{hDiv}, rather than derived from a basis of
periodic tensors, and the \emph{Lem1t} almost-orthonormal-implies-independent
lemma is not used. Restoring the spanning clause and routing through
\emph{Lem1t} was classified under
\href{https://github.com/LionSR/TNLean/issues/81}{issue \#81}. The independence
half is formalized; no unrestricted spanning assertion is presently attributed
to the source consequence.

\section{Normalized multiplicity theorem and remaining transport restriction}

\textbf{Source.}\footnote{\cite{DeLasCuevas2017Irreducible}, equations
\emph{bdnr} and \emph{Bbdnr}, lines 286--305 and 575--585; the normalization
reduction, lines 313--332; and theorem \emph{thm:bd}, lines 613--632.} The
irreducible forms used in the proportional Fundamental Theorem are
\begin{align}
  A^i&=\bigoplus_j (R_j\otimes A_j^i),&
  B^i&=\bigoplus_k (S_k\otimes B_k^i),
\end{align}
where the diagonal matrices $R_j$ and $S_k$ retain the multiplicities and
weights of repeated copies of each periodic basis tensor.

\textbf{Multiplicity-bearing normalized theorem.} The declaration
\path{PeriodicOverlapHypothesis.ofSectorDecompositions} uses literal sector
decompositions
\begin{align}
  P^i&=\bigoplus_{j,q}(\mu_{j,q}A_j^i),&
  Q^i&=\bigoplus_{k,r}(\nu_{k,r}B_k^i),
\end{align}
where every complex weight is nonzero.  This flattened presentation is the
coordinate form of $\bigoplus_j(R_j\otimes A_j^i)$ and
$\bigoplus_k(S_k\otimes B_k^i)$, with
\begin{align}
  \operatorname{tr}(R_j^N)&=\sum_q\mu_{j,q}^N,&
  \operatorname{tr}(S_k^N)&=\sum_r\nu_{k,r}^N.
\end{align}
It assumes proportionality only at positive lengths, permits different total
bond dimensions, and allows the proportionality scalar to depend freely on
the length.  Thus the multiplicity, literal-representation, and empty-word
restrictions of the earlier scalar theorem are absent.

The proof makes the paper's terse contradiction at line~631 explicit.  It
partitions one basis by mixed-overlap decay, passes to a common-period
subsequence, and compares the grouped coefficients in an eventually linearly
independent coproduct-indexed family.  Linear independence forces
$\operatorname{tr}(R_{j_0}^{pN})=0$ for every sufficiently large $N$.
Geometric extrapolation applied to the nonzero numbers $\mu_{j_0,q}^p$ then
forces the same power sum to vanish at exponent zero, contradicting the
positive multiplicity of the sector.  In particular, the proof never treats
the repeated copies of $A_{j_0}$ as independent vectors.

The earlier declaration
\path{PeriodicOverlapHypothesis.ofIsIrreducibleForm} remains as the
one-weight-per-listed-block specialization.  Its hypothesis still asserts
matrix-product-vector equality at the empty word and may describe an ambient
tensor only through such an equality.  These are restrictions of that older
formulation, not of the new literal multiplicity-bearing theorem.

\textbf{Pure normalization theorem.} The predicate
\path{IsSpectrallyPeriodic} records the source assumptions on a basis block:
irreducibility, spectral radius one, positive period, and the prescribed
roots-of-unity peripheral spectrum.  The declaration
\path{IsSpectrallyPeriodic.exists_isPeriodic_tpGauge} applies the positive
adjoint Perron eigenmatrix to obtain a left-canonical periodic block by pure
similarity, exactly as in
\cite[lines~313--332]{DeLasCuevas2017Irreducible}.  The Perron eigenvalue is
forced to equal one, so no scalar rescales the block.  The peripheral spectrum
is preserved by similarity.  Thus this normalization leaves every
matrix-product vector unchanged and introduces no extra power into the
multiplicity coefficients.

\textbf{Remaining sectorwise transport restriction.} The normalized overlap
theorem still takes a literal sector decomposition whose basis blocks already
satisfy \path{IsPeriodic}.  The formal development does not yet replace every
basis block by its Perron representative inside a new sector decomposition,
assemble the individual similarities into a block-diagonal gauge, or transport
the overlap hypothesis back to the original basis.  Applying the theorem to an
arbitrary original tensor through a lower-dimensional representative also
requires a positive-length transport statement of the kind allowed at
lines~266--271.

The declarations \leanid{fundamentalTheorem_periodic_proportional} and
\leanid{fundamentalTheorem_periodic_proportional_of_isPeriodic} begin after
this step: they accept the periodic-overlap hypothesis, or its non-decaying
partner fields, as premises and prove the resulting finite block matching.
Together with the new normalized overlap theorem they give the block matching for
literal normalized irreducible forms.  The source theorem itself remains
unformalized until the preceding sectorwise transport step is proved;
the exact source-labelled blueprint statement must therefore remain
\emph{not ready}.

\textbf{Elimination plan.} Apply the pure Perron normalization to every basis
block and assemble the block gauges while leaving all multiplicity entries
unchanged.  Prove invariance of proportional positive-length MPVs, overlap
decay, pairwise non-repetition, and the final repeated-block relation under
these similarities.  A separate positive-length representative
transport then covers the source-allowed lower-dimensional representative.
The remaining sectorwise transport work is tracked by
\href{https://github.com/LionSR/TNLean/issues/5342}{issue \#5342}; the subsequent
multiplicity-power extraction is tracked by
\href{https://github.com/LionSR/TNLean/issues/5344}{issue \#5344}.

\begin{thebibliography}{9}
\bibitem{DeLasCuevas2017Irreducible}
G.~de las Cuevas, J.~I.~Cirac, N.~Schuch, and D.~P\'erez-Garc\'ia,
\emph{Irreducible forms of matrix product states: Theory and applications},
J.~Math.~Phys.\ \textbf{58}, 121901 (2017), arXiv:1708.00029.

\bibitem{CiracPerezGarcia2017Matrix}
J.~I.~Cirac, D.~P\'erez-Garc\'ia, N.~Schuch, and F.~Verstraete,
\emph{Matrix product states and projected entangled pair states: Concepts,
symmetries, theorems}, Ann.\ Phys.\ \textbf{378}, 100 (2017)
(``Ci15'' in arXiv:1708.00029).

\bibitem{Wolf2012Quantum}
M.~M.~Wolf, \emph{Quantum Channels and Operations: Guided Tour}, unpublished
lecture notes (2012).
\end{thebibliography}

\end{document}
